\chapter{Rangs, fonctions et origines}\label{ch:rangs}

\section{Le clarissimat des défunts}

La première version du catalogue réunissait sous la rubrique « \emph{clarissimus} » la titulature des consuls, \emph{v(iri) c(larissimi)}, et le rang des personnes commémorées. Les deux sont désormais séparés. La titulature apparaît dans environ quatre-vingt-dix textes de l’ensemble C, presque toujours dans la formule consulaire. Le rang sénatorial d’une personne commémorée ou d’un commanditaire n’est attesté que dans onze textes : \nnums{15}, 19, 36, 44, 179, 213, 262, 306, 330, 430 et 452.

Ces onze textes forment un groupe hétérogène. Cassia Pisonis, \emph{c(larissima) f(emina)}, dont l’épitaphe est gravée sur un sarcophage, est née en 316 et morte en 346 (\nnum{15}). Au \nnum{179}, une même dalle commémore Simplicius, \emph{v(ir) c(larissimus)}, Didyme et Roscia Calchedonia, \emph{c(larissimae) f(eminae)}, et Innocentius, \emph{[dormi]entes in pace} (participe en partie restitué). Au \nnum{213}, Eustolius, \emph{v(ir) c(larissimus)}, repose avec sa fille. Egnatia Susanna, dite aussi Abita, est une \emph{c(larissima) p(uella)} morte à cinq ans (\nnum{330}). Les autres textes sont fragmentaires (\nnums{44}, 262, 430, 452).

Deux textes posent des questions d’identification.

\emph{Le \nnum{36}.} Clodius Adelfius, \emph{v(ir) c(larissimus) ex praefectis urbis}, fait graver l’épitaphe de son « épouse incomparable » et la sienne. La bibliographie l’identifie habituellement à Clodius Celsinus Adelphius, préfet de la Ville en 351 et époux de la poétesse chrétienne Faltonia Betitia Proba. L’identification est vraisemblable mais n’a pas été contrôlée sur la \emph{Prosopography of the Later Roman Empire}. Elle est en tout cas la seule source de l’année : le texte ne contient aucune date, et l’année 351 de la base est un terme postérieur, non l’année du décès (chapitre~\ref{ch:fondement}).

\emph{Le \nnum{306}.} Lucceia, \emph{c(larissima) f(emina)}, est dite fille de Viventius, \emph{c(larissimae) r(ecordationis) v(ir) ex praefecto praetorio et urbis aeternae}. L’identification de ce Viventius au préfet de la Ville de la fin des années 360 relève de la même vérification prosopographique. Lucceia fait ensevelir dans son propre sarcophage la vierge Maximilla, originaire de Pannonie, « par amitié », avec la mère de celle-ci, Nunita, \emph{matrona diaconi} (chapitre~\ref{ch:textes_decisifs}).

\section{Clercs, fonctions et métiers}

Le tableau~\ref{tab:fonctions} réunit les mentions de fonctions ecclésiastiques, de charges et de métiers. Le relevé porte sur l’ensemble B, qui comprend les actes d’achat, et sur la colonne de Saint-Paul (\nnum{321}). L’astérisque signale une attestation dont le mot est entièrement restitué.

% généré par scripts/04_generated.py
\begin{table}[tbp]\centering\small\caption{Fonctions, rangs et métiers mentionnés (ensemble B et \nnum{321} ; l’astérisque signale une attestation seulement restituée). Le rang n’est pas toujours celui du défunt.}\label{tab:fonctions}
\begin{tabular}{@{}l r p{0.5\linewidth}@{}}\toprule Terme & Textes & Notices\\\midrule
\lat{lector} & 4 & 67*, 202, 269, 421 \\
\lat{presbyter} & 8 & 93*, 105, 134, 289*, 310, 378, 387*, 428* \\
\lat{diaconus} & 3 & 134, 306, 378 \\
\lat{episcopus} & 4 & 321, 378, 387, 400 \\
\lat{fossor} & 4 & 209, 217, 424, 478 \\
\lat{medicus} & 1 & 296 \\
\lat{pictor} & 1 & 229* \\
\lat{numerarius} & 1 & 429 \\
\lat{scholasticus} & 1 & 469 \\
\lat{quadratarius} & 1 & 189* \\
\lat{armararius} & 1 & 365 \\
\lat{praefectus} & 3 & 36, 72, 306 \\
\lat{praeses (ex praesidibus)} & 1 & 20 \\
\lat{tribunus} & 1 & 321* \\
\lat{sacerdos} & 1 & 289* \\
\lat{matrona (diaconi, presbyteri)} & 3 & 306, 378*, 428 \\
\bottomrule\end{tabular}\end{table}


\emph{Les lecteurs.} Trois lecteurs sont rattachés à une église de la Ville : Cinnamius Opas, \emph{lector tituli Fasciolae}, mort à 46 ans en 377 (\nnum{202}) ; Leopardus, \emph{lector de Pudentiana}, mort à 24 ans en 384 (\nnum{269}) ; Pascentius, \emph{lector de Fasc[iola]}, mort à 21 ans en 398 (\nnum{421}). Un quatrième, Iulius Innocentius, est dit \emph{lect[or] ec(c)les[iae sanctae c]at(h)olic(a)e} et serait mort à 73 ans en 362 (\nnum{67}) ; mais le mot \emph{lector} est en partie restitué, et l’expression \emph{ecclesiae sanctae catholicae} l’est presque entièrement. Ces âges montrent que la fonction de lecteur n’est pas réservée aux jeunes gens. Cinnamius Opas est en outre qualifié d’\emph{amicus pauperum}.

\emph{Prêtres, diacres et évêques.} Le \nnum{134} commémore un prêtre, Dulcitius, et un diacre, Gregorius, frères, avec leur sœur. Le \nnum{310} est l’épitaphe qu’un prêtre, Gaudentius, fait graver pour lui-même et pour son épouse Severa, \emph{casta}. Au \nnum{428}, Remigia, \emph{matrona} du prêtre Laurentius, fait graver son épitaphe ; mais le titre de Laurentius n’est conservé que par les lettres \emph{[p]r(es)b(yter)}, et la mention de l’\emph{ecclesia catholica} est restituée. Au \nnum{306}, Nunita est \emph{matrona diaconi} : l’EDCS supprime le S final gravé (\emph{diaconi\{s\}}), et la lecture « épouse d’un diacre » dépend de cette correction. Ces textes documentent l’ancrage familial du clergé romain et le mariage des clercs à la fin du \textsc{iv}\textsuperscript{e} siècle.

Le \nnum{378} développe une généalogie cléricale : à la suite du père, le prêtre Limenius, le texte énumère une série d’ascendants et de parents (\emph{nepos}, \emph{pronepos}, \emph{abnepos}, \emph{socrus}) parmi lesquels un diacre de la cinquième région, un évêque Marcellus, une \emph{matrona presbyteri} et un évêque Ursinus d’Albano. La structure exacte de cette généalogie dépend de nombreuses restitutions. Surtout, elle n’est transmise que par l’EDCS : aucune des fiches de l’EDB rattachées au monument ne la contient (chapitre~\ref{ch:evenements}). Elle doit être contrôlée sur l’ICVR III 8345 avant toute exploitation prosopographique.

Trois mentions d’évêques concernent l’évêque de Rome ou un évêque dont l’identité est discutée. Sur la colonne de Saint-Paul, \emph{Siricius episcopus tota mente devotus} (\nnum{321}, non funéraire) atteste l’intervention de l’évêque dans la construction de la basilique. Au \nnum{387}, la formule \emph{sub episcopatu Si[rici]} est en grande partie restituée, de même que l’emplacement \emph{[ad limina Petri et] Pauli}. Au \nnum{400}, l’emplacement a été acheté par Viricunda, épouse du défunt, \emph{Pascasio ep(is)c(opo)} ; l’EDB reprend la note de Diehl selon laquelle aucun évêque de Rome ne porte ce nom. La fonction de ce Pascasius dans la transaction reste indéterminée (chapitre~\ref{ch:sepultures}).

\emph{Charges et métiers.} Le \nnum{20} commémore Flavius Gaianus, \emph{ex praesidibus}, ancien gouverneur de province. Le \nnum{429} nomme Martielis, \emph{numerarius vicariae sedis urbis Romae aeternae}, comptable de l’administration du vicaire de la Ville. Le \nnum{469} mentionne deux \emph{scholastici}. Le \nnum{296} est l’épitaphe de Rapetiga, médecin, \emph{civis Hispanus}, mort à environ 25 ans. Un peintre (\emph{pict[or]}, \nnum{229}) et un \emph{[qua]dratarius} (\nnum{189}) sont attestés par des mots en partie restitués. Les \emph{fossores} apparaissent uniquement comme vendeurs d’emplacements (\nnums{209}, 217, 424 et 478 ; chapitre~\ref{ch:sepultures}).

\section{Dater par le préfet de la Ville}

Le \nnum{72} associe à la formule consulaire de 362 (\emph{Mamertino et Nevitta conss.}) la mention du préfet de la Ville, \emph{pr(a)ef(ecto) urb(i) Maximo}. Cette double datation, exceptionnelle dans le corpus, appartient à un texte qui commence par \emph{D(is) M(anibus) s(acrum)} et se poursuit par \emph{requievit in pace} : la formule païenne traditionnelle et la formule chrétienne coexistent sur la même plaque (chapitre~\ref{ch:mort_esperance}).

\section{Origines déclarées}

Quelques défunts sont présentés par leur origine : Quirillus, \emph{civis Armeniacus Cappadox} (\nnum{275}) ; Rapetiga, \emph{civis Hispanus} (\nnum{296}) ; Maximilla, \emph{civis Pannonia} (\nnum{306}, ligne dont la présence doit être contrôlée sur l’édition, car elle manque dans le texte de l’EDB) ; Vitalia, \emph{εὐγενεστάτη Κωνσταντινουπολίτισσα} (\nnum{168}, en partie restitué) ; Alexandros, du village de Mikra Kômè (\nnum{348}) ; selon l’EDB, un originaire de Théodosiopolis (\nnum{447}, restitué). Au \nnum{19}, Flavius Aurelius Casarius, dit Santulus, est venu de Préneste à Rome, où il a vécu vingt-deux jours ; l’indication \emph{[ante annos d]uos} est restituée. Le \nnum{183} situe la défunte, Petronia, \emph{virgo}, dans la neuvième région de la Ville, \emph{a lacu cunic(u)li} ; la localisation de ce repère relève de la topographie romaine et n’est pas proposée ici.

Ces mentions sont rares et ne permettent aucune estimation de la part des immigrés dans la population chrétienne. Elles montrent seulement que l’origine est parfois jugée digne d’être gravée, surtout pour des personnes venues de loin.
